\documentclass[11pt]{article}
\usepackage[margin=0.88in]{geometry}
\usepackage{amsmath,amssymb,lmodern}
\usepackage[T1]{fontenc}
\usepackage{needspace}
\usepackage[colorlinks=true,linkcolor=blue,urlcolor=blue]{hyperref}
\setlength{\emergencystretch}{3em}
\setlength{\parskip}{0.35em}
\setcounter{secnumdepth}{0}
\providecommand{\tightlist}{\setlength{\itemsep}{0pt}\setlength{\parskip}{0pt}}
\title{Euclidean persistence and sharp stability for anisotropic eigenvalue--torsion products}
\author{Henry Zweiman}
\date{September 15, 2026}
\begin{document}
\maketitle
\section{Abstract}\label{abstract}

We optimize the product of the first Dirichlet eigenvalue and a power of
torsional rigidity on a fixed Euclidean ball, allowing the anisotropic
energy to vary over all normalized symmetric seminorms. For every
dimension \(n\geq2\) and every energy exponent \(1<p<\infty\), we prove
that the Euclidean norm uniquely maximizes the product for all
sufficiently small positive torsion exponents and uniquely minimizes it
for all sufficiently large exponents. In either regime, the product
deficit is equivalent to the spherical average of \(1-H^p\).
Consequently it controls the supremum norm distance from the Euclidean
norm with the optimal power \((n+1)/2\). The proof combines radial
variational tests, convex symmetrization, and an integral comparison
between support and radial deficits of convex bodies inside a ball. For
quadratic growth, this answers the ball cases of two
Euclidean-optimality questions of Buttazzo and Fernandes Horta, beyond
the quadratic-seminorm class. No regularity or strict convexity of the
competing anisotropies is assumed.

\textbf{Keywords:} anisotropic spectral optimization; torsional
rigidity; nonlinear eigenvalue; convex symmetrization; quantitative
stability.

\section{1. Introduction and main
results}\label{introduction-and-main-results}

Fix the Euclidean unit ball \(B\subset\mathbb R^n\), where \(n\geq2\),
and an exponent \(1<p<\infty\). Denote its volume by \(V\), write
\(S=\partial B\), and let \(\sigma\) be probability surface measure on
\(S\). We consider the class \(\mathcal H\) of symmetric seminorms
\(H:\mathbb R^n\to[0,\infty)\) normalized by

\[
\max_{\theta\in S}H(\theta)=1.
\tag{1.1}
\]

Thus \(H\) is convex, satisfies \(H(t\xi)=|t|H(\xi)\), and may vanish on
a nonzero subspace. The Euclidean norm is denoted by \(E(\xi)=|\xi|\).
For \(H\in\mathcal H\) define

\[
\begin{aligned}
\lambda_H&=\inf_{u\in W^{1,p}_0(B)\setminus\{0\}}
\frac{\int_B H(\nabla u)^p\,dx}{\int_B|u|^p\,dx},\\
T_H&=\left(
\sup_{u\in W^{1,p}_0(B)\setminus\{0\}}
\frac{|\int_Bu\,dx|^p}{\int_BH(\nabla u)^p\,dx}
\right)^{1/(p-1)}.
\end{aligned}
\tag{1.2}
\]

The dependence on \(n,p\) is suppressed. These quantities are finite and
strictly positive even for degenerate seminorms; a direct proof is given
below. The definition does not require attainment in the degenerate
case. Our torsion convention agrees with the integral of the solution of
the torsion equation in the usual coercive setting. In particular,
\(T_H^{p-1}\) is the variational quotient before taking the root in
(1.2).

The normalized product and its angular deficit are

\[
R_q(H)=\frac{\lambda_H T_H^q}{\lambda_E T_E^q},
\qquad
\delta_p(H)=\int_S\bigl(1-H(\theta)^p\bigr)\,d\sigma(\theta).
\tag{1.3}
\]

The exponent \(q>0\) is independent of the energy exponent \(p\). Since
\(H\leq E\), decreasing the anisotropy lowers the eigenvalue and raises
the torsion. Separate monotonicity statements therefore do not determine
the product. Even convergence of optimizers to \(E\) as \(q\) tends to
an endpoint would not establish that \(E\) is the exact optimizer at a
fixed finite \(q\).

\Needspace{6\baselineskip}\textbf{Theorem 1.1 (Euclidean persistence and deficit equivalence).}
For every \(n\geq2\) and \(1<p<\infty\), there are numbers

\[
0<q_-(n,p)<p-1<q_+(n,p)<\infty
\tag{1.4}
\]

with the following properties. The Euclidean norm uniquely maximizes
\(\lambda_H T_H^q\) over \(\mathcal H\) whenever \(0<q\leq q_-\), and
uniquely minimizes it whenever \(q\geq q_+\). For each fixed \(q\) in
these regimes, set

\[
D_q(H)=
\begin{cases}
1-R_q(H),&0<q\leq q_-,\\
R_q(H)-1,&q\geq q_+.
\end{cases}
\tag{1.5}
\]

There are constants \(0<a_{n,p,q}\leq b_{n,p,q}<\infty\) such that

\[
a_{n,p,q}\delta_p(H)\leq D_q(H)
\leq b_{n,p,q}\delta_p(H)
\qquad(H\in\mathcal H).
\tag{1.6}
\]

The theorem applies to the full infinite-dimensional class of normalized
symmetric seminorms, with no ellipticity constant prescribed in advance.
In particular it includes polyhedral norms and seminorms of every
positive rank. Explicit, though nonoptimal, choices of the thresholds
are provided below.

The associated uniform distance is

\[
\varepsilon(H)=\|H-E\|_{L^\infty(S)}
=\max_{\theta\in S}\bigl(1-H(\theta)\bigr).
\tag{1.7}
\]

\Needspace{6\baselineskip}\textbf{Theorem 1.2 (optimal stability power).} Under the assumptions
and exponent restrictions of Theorem 1.1, there is \(c_{n,p,q}>0\) such
that

\[
D_q(H)\geq c_{n,p,q}\varepsilon(H)^{(n+1)/2}
\qquad(H\in\mathcal H).
\tag{1.8}
\]

The power \((n+1)/2\) cannot be replaced by a smaller one, even if the
competitors are restricted to genuine norms. More precisely, there are
norms \(H_\varepsilon\in\mathcal H\), for \(0<\varepsilon<1/2\), with
\(\varepsilon(H_\varepsilon)=\varepsilon\) and

\[
D_q(H_\varepsilon)\leq C_{n,p,q}\varepsilon^{(n+1)/2}.
\tag{1.9}
\]

By scaling the spatial variable, the same normalized conclusions hold on
any Euclidean ball. The thresholds do not depend on its radius: the
eigenvalue scales by the inverse \(p\)-th power of the radius and
torsion by the power \(n+p/(p-1)\), and these factors cancel in \(R_q\).

\subsubsection{Relation to previous
work}\label{relation-to-previous-work}

Buttazzo and Fernandes Horta {[}BFH{]} study eigenvalue-torsion products
while varying normalized seminorms on a fixed domain. Their Section 6
asks whether the norm optimizers in the small-\(q\) maximum and
large-\(q\) minimum regimes are Euclidean. For balls, their Theorem 4.10
proves maximality for \(q\leq1\) within the class of quadratic
seminorms. Theorem 1.1 answers both questions on balls in the full
symmetric-seminorm class, and treats the general \(p\)-growth setting
also proposed in their final remarks. It does not resolve their
questions for arbitrary fixed domains or determine the optimal
thresholds.

The optimization of the eigenvalue alone, including rigidity of
Euclidean maximality, was studied by Fernandes Horta and Montenegro
{[}HM{]}; see their Theorem 1.5. We do not claim that endpoint result as
new. Our spectral lower estimate uses the convex symmetrization theorem
of Alvino, Ferone, Trombetti and Lions {[}AFTL{]}. The body and
normalization conventions for this input are stated explicitly in
Section 4. The convex-geometric part uses support and radial deficits,
quantities related to the intrinsic and dual volume deviations studied
in {[}BHK{]}. We provide the geometric estimate needed here in full and
make no separate priority claim for the elementary cap geometry.

The main analytic point is a common linear control, in \(\delta_p(H)\),
of the spectral loss and torsion gain near \(E\). Directional first
variation alone gives no such uniform estimate over increasingly
concentrated angular perturbations. The support-to-radial comparison in
Section 3 supplies this missing control without smoothness assumptions.
It also produces the sharp supremum norm stability power.

\section{2. Convex representation and global variational
bounds}\label{convex-representation-and-global-variational-bounds}

Every \(H\in\mathcal H\) is the support function of the compact
centrally symmetric convex set

\[
K=\{x\in\mathbb R^n:x\cdot\xi\leq H(\xi)
\text{ for every }\xi\in\mathbb R^n\},
\qquad H=h_K.
\tag{2.1}
\]

Indeed, the support representation follows from separation for finite
continuous sublinear functions. Normalization gives
\(K\subset\overline B\) and \(K\cap S\ne\varnothing\). To check the
latter, the maximum of \(h_K\) on \(S\) equals the maximum of \(|x|\)
over \(K\). In particular, there is \(v\in K\cap S\), and symmetry
implies

\[
[-v,v]\subset K,
\qquad H(\xi)\geq|v\cdot\xi|.
\tag{2.2}
\]

Also \(|H(\xi)-H(\eta)|\leq|\xi-\eta|\). When \(K\) has interior, write
\(\rho_K\) for its radial function and \(H^\circ\) for the dual norm of
\(H\). Then

\[
K=\{H^\circ\leq1\},\qquad
H^\circ(\theta)=\rho_K(\theta)^{-1}.
\tag{2.3}
\]

Thus \(K\) is the polar of the unit ball of \(H\), not the unit ball of
\(H\) itself.

\Needspace{6\baselineskip}\textbf{Lemma 2.1 (global bounds).} Put

\[
m_H=\int_SH^p\,d\sigma=1-\delta_p(H),
\qquad p'=\frac p{p-1}.
\tag{2.4}
\]

For every \(H\in\mathcal H\),

\[
2^{-p}\leq\lambda_H\leq m_H\lambda_E,
\qquad
m_H^{-1/(p-1)}T_E\leq T_H\leq V2^{p'}.
\tag{2.5}
\]

Moreover, \(m_H\geq\int_S|\theta\cdot v|^p\,d\sigma>0\), where this last
integral depends only on \(n,p\).

\emph{Proof.} For \(f\in C_c^\infty(I)\) on an interval of length
\(L\leq2\), the fundamental theorem of calculus and
Hölder\textquotesingle s inequality imply

\[
\int_I|f|^p\leq L^p\int_I|f'|^p
\leq2^p\int_I|f'|^p.
\tag{2.6}
\]

Apply this on chords parallel to \(v\) and integrate over the orthogonal
variables. Using (2.2) and density in \(W^{1,p}_0(B)\) gives

\[
\int_B|u|^p\leq2^p\int_BH(\nabla u)^p.
\tag{2.7}
\]

This proves the eigenvalue lower bound. Combining (2.7) with
\(|\int_Bu|^p\leq V^{p-1}\int_B|u|^p\) proves the torsion upper bound in
(2.5). In particular, none of the denominators in (1.2) vanish for
nonzero admissible functions.

A Euclidean first eigenfunction can be chosen radial and nonnegative by
Schwarz symmetrization of a minimizer. Existence follows by the direct
method and the compact embedding into \(L^p(B)\). For every radial
Sobolev function \(f\), homogeneity and polar coordinates give

\[
\int_BH(\nabla f)^p=m_H\int_B|\nabla f|^p.
\tag{2.8}
\]

Testing the eigenvalue with a radial Euclidean minimizer proves its
upper bound. The Euclidean torsion function and its integral are

\[
w_E(x)=\frac{n^{-1/(p-1)}}{p'}(1-|x|^{p'}),
\qquad
T_E=\frac{V n^{-1/(p-1)}}{n+p'}.
\tag{2.9}
\]

In fact \(|\nabla w_E|^{p-2}\nabla w_E=-x/n\), so \(-\Delta_pw_E=1\) in
distributions. Testing this identity with \(w_E\) gives

\[
\int_B|\nabla w_E|^p=\int_Bw_E=T_E.
\tag{2.10}
\]

Hölder applied to the same weak identity shows that \(w_E\) attains the
Euclidean torsion quotient. Equations (2.8)-(2.10) then give its lower
bound in (2.5). The final assertion follows by integrating (2.2) on
\(S\). \(\square\)

We shall use the fixed constants

\[
\ell=\min\left\{\frac12,\frac{2^{-p}}{\lambda_E}\right\},
\qquad
M=\max\left\{2,\frac{V2^{p'}}{T_E}\right\}.
\tag{2.11}
\]

They give the global bounds \(\lambda_H/\lambda_E\geq\ell\) and
\(T_H/T_E\leq M\). Their sharp values are unnecessary.

\section{3. An integral comparison of support and radial
deficits}\label{an-integral-comparison-of-support-and-radial-deficits}

For \(\nu\in S\) and \(r>0\), let
\(C(\nu,r)=\{\theta\in S:|\theta-\nu|<r\}\) be a chordal cap. The
notation \(aC\) denotes the cap with the same center and radius
multiplied by \(a\).

\Needspace{6\baselineskip}\textbf{Lemma 3.1 (support-to-radial comparison).} There is
\(A_n<\infty\) such that every compact convex set \(K\) with
\(\tfrac12\overline B\subset K\subset\overline B\) satisfies

\[
\int_S(1-\rho_K)\,d\sigma
\leq A_n\int_S(1-h_K)\,d\sigma.
\tag{3.1}
\]

Symmetry is not required in this lemma.

\emph{Proof.} We first record a cap estimate. Suppose that
\(h_K(\nu)=1-e\), where \(0<e\leq1\). Since
\(K\subset\{x\in\overline B:x\cdot\nu\leq1-e\}\), its support function
at a direction \(\eta\) making angle \(\alpha\leq\sqrt e/4\) with
\(\nu\) satisfies

\[
\begin{aligned}
h_K(\eta)
&\leq(1-e)\cos\alpha+\sqrt{2e-e^2}\sin\alpha\\
&\leq1-e+\frac{\sqrt2}{4}e
\leq1-\frac e2.
\end{aligned}
\tag{3.2}
\]

For completeness, \(\cos\alpha>1-e\) because
\(1-\cos\alpha\leq\alpha^2/2\leq e/32\). Thus the unconstrained support
point \(\eta\) is excluded by the cutting plane, and maximizing the
linear functional on the truncated ball puts the maximizing point on
that plane. Decomposition into the \(\nu\) and \(\nu^\perp\) components
gives the first line of (3.2). The last estimates use \(1-e\geq0\) and
\(\sin\alpha\leq\alpha\).

Set \(r=\sqrt e/8\). On \(C(\nu,r)\) the angle is less than
\(2r\leq\sqrt e/4\), so this cap lies in \(\{1-h_K\geq e/2\}\).

Now fix \(\theta\in S\setminus K\). Let \(z\) be its metric projection
onto \(K\), and put

\[
d=|\theta-z|,\qquad
\nu=\frac{\theta-z}{d},\qquad
e=1-h_K(\nu).
\tag{3.3}
\]

The projection inequality gives \(z\cdot\nu=h_K(\nu)\), hence
\(\theta\cdot\nu=1-e+d\). Therefore

\[
0<d\leq e\leq\frac12,
\qquad
|\theta-\nu|^2=2(e-d)\leq2e.
\tag{3.4}
\]

The cap \(C=C(\nu,\sqrt e/8)\) from (3.2) consequently has
\(\theta\in12C\).

Fix \(t>0\) and write
\(E_t=\{\theta\in S:\operatorname{dist}(\theta,K)>t\}\). For each
\(\theta\in E_t\) choose the cap just constructed; its associated \(e\)
exceeds \(t\). Every compact subset \(L\subset E_t\) is covered by
finitely many of the open caps \(12C\). Order the corresponding original
caps by decreasing radius, and retain a cap precisely when it is
disjoint from every previously retained cap. The retained caps \(C_i\)
are disjoint. Any discarded cap of radius \(r\) intersects a retained
cap of radius \(r_i\geq r\), so its twelvefold enlargement lies in
\(14C_i\), and hence in \(30C_i\). Thus \(L\subset\bigcup_i30C_i\).

Spherical cap measure is doubling: there is \(D_n<\infty\) such that

\[
\sigma(C(\nu,30r))\leq D_n\sigma(C(\nu,r))
\quad\left(0<r\leq\frac{1}{8\sqrt2}\right).
\tag{3.5}
\]

This follows directly by writing cap area in polar coordinates. The
ratio tends to \(30^{n-1}\) as \(r\downarrow0\) and is bounded on every
remaining compact interval; a cap whose chordal radius exceeds \(2\) is
the entire sphere. Each retained cap is contained in \(\{1-h_K>t/2\}\)
by (3.2). Disjointness and (3.5) give

\[
\sigma(L)\leq D_n\sum_i\sigma(C_i)
\leq D_n\sigma\{1-h_K>t/2\}.
\tag{3.6}
\]

By inner regularity this holds with \(L=E_t\). Integrating in \(t\)
yields

\[
\int_S\operatorname{dist}(\theta,K)\,d\sigma
\leq2D_n\int_S(1-h_K)\,d\sigma.
\tag{3.7}
\]

The gauge \(g_K\) is \(2\)-Lipschitz since \(K\) contains
\(\tfrac12\overline B\). Indeed, subadditivity gives
\(g_K(x)-g_K(y)\leq g_K(x-y)\leq2|x-y|\), and one can interchange
\(x,y\). For the projection \(z\) above, \(g_K(z)\leq1\), so

\[
\rho_K(\theta)^{-1}-1
\leq g_K(\theta)-g_K(z)
\leq2\operatorname{dist}(\theta,K).
\tag{3.8}
\]

Since \(0<\rho_K\leq1\), we have \(1-\rho_K\leq\rho_K^{-1}-1\).
Equations (3.7)-(3.8) prove (3.1) with \(A_n=4D_n\). \(\square\)

The estimate transfers angular support loss to the radial moments needed
by both variational problems. A supremum norm comparison alone would not
give the linear dependence on \(\delta_p\) used below.

\section{4. Uniform estimates near the Euclidean
norm}\label{uniform-estimates-near-the-euclidean-norm}

Set

\[
\eta_0=\frac12\bigl(1-(3/4)^p\bigr)\sigma(C(\nu,1/4)),
\tag{4.1}
\]

which is positive and independent of \(\nu\). If \(\min_SH\leq1/2\), the
\(1\)-Lipschitz property implies \(H\leq3/4\) on a cap of radius
\(1/4\). Thus \(\delta_p(H)\geq2\eta_0\). In particular, if
\(\delta_p(H)\leq\eta_0\), then \(H>1/2\) on \(S\) and \(K\) contains
\(\tfrac12\overline B\).

For the remainder of this section write
\(\delta=\delta_p(H)\leq\eta_0\). Since \(1-H\leq1-H^p\) on \(S\), Lemma
3.1 gives

\[
\int_S(1-\rho_K)\,d\sigma\leq A_n\delta.
\tag{4.2}
\]

Polar integration and \(1-r^n\leq n(1-r)\) for \(0\leq r\leq1\) yield

\[
\frac{|K|}{V}=\int_S\rho_K^n\,d\sigma
\geq1-nA_n\delta.
\tag{4.3}
\]

Similarly, on \([1/2,1]\) the derivative of \(r^{-p'}\) has magnitude at
most \(p'2^{p'+1}\), and therefore

\[
\int_S(H^\circ)^{p'}\,d\sigma
\leq1+p'2^{p'+1}A_n\delta.
\tag{4.4}
\]

Fix

\[
C=\max\{nA_n,p'2^{p'+1}A_n,1\},
\qquad
\eta=\min\{\eta_0,(2C)^{-1},1/4\}.
\tag{4.5}
\]

\Needspace{12\baselineskip}\textbf{Lemma 4.1 (spectral and torsion control).} If
\(\delta_p(H)\leq\eta\), then

\[
(1-C\delta_p(H))^{p/n}
\leq\frac{\lambda_H}{\lambda_E}\leq1-\delta_p(H),
\tag{4.6}
\]

and

\[
(1-\delta_p(H))^{-1/(p-1)}
\leq\frac{T_H}{T_E}\leq1+C\delta_p(H).
\tag{4.7}
\]

\emph{Proof.} Only the lower spectral and upper torsion bounds remain.
Convex symmetrization {[}AFTL, Theorem 3.1 and equation (3.5){]}
implies, for \(u\in W^{1,p}_0(B)\) and its Euclidean Schwarz
rearrangement \(u^\#\),

\[
\int_BH(\nabla u)^p\,dx
\geq\left(\frac{|K|}{V}\right)^{p/n}
\int_B|\nabla u^\#|^p\,dx.
\tag{4.8}
\]

Here one first replaces \(u\) by \(|u|\); evenness of \(H\) preserves
the energy. To specify the normalization in (4.8), let
\(\mu(t)=|\{|u|>t\}|\). The energy of the convex rearrangement has the
level-set expression

\[
\int_0^\infty
\bigl[n|K|^{1/n}\mu(t)^{1-1/n}\bigr]^p
(-\mu'(t))^{1-p}\,dt,
\tag{4.9}
\]

in the usual rearrangement interpretation. The corresponding expression
for \(u^\#\) replaces \(|K|\) by \(V\). The convex rearrangement theorem
decreases the original energy; comparison of these expressions gives
(4.8), with the general Sobolev statement supplied by that theorem. The
body denoted \(K\) in Section 2 of {[}AFTL{]} is the unit ball of \(H\);
its polar is the body called \(K\) here. Their auxiliary volume
normalization can be removed by scaling \(H\): both sides of (4.8) scale
by its \(p\)-th power. Their hypotheses allow nonsmooth convex norms
with positive coercivity, which holds in the present neighborhood. No
regularity approximation of \(H\) or equality case is needed.

Equimeasurability preserves the \(L^p\) denominator in the Rayleigh
quotient. Taking the infimum in (4.8) gives
\(\lambda_H/\lambda_E\geq(|K|/V)^{p/n}\), and (4.3) proves (4.6).

For torsion use the Euclidean flux

\[
F=|\nabla w_E|^{p-2}\nabla w_E=-\frac xn.
\tag{4.10}
\]

For every \(u\in W^{1,p}_0(B)\), integration by parts, duality of norms
and Hölder\textquotesingle s inequality give

\[
\left|\int_Bu\right|^p
=\left|\int_BF\cdot\nabla u\right|^p
\leq\left(\int_BH^\circ(F)^{p'}\right)^{p-1}
\int_BH(\nabla u)^p.
\tag{4.11}
\]

Take the supremum and the power \(1/(p-1)\) in (1.2). Radiality of \(F\)
and (2.10) imply

\[
T_H\leq\int_BH^\circ(F)^{p'}
=T_E\int_S(H^\circ)^{p'}\,d\sigma.
\tag{4.12}
\]

Equation (4.4) proves the upper torsion bound. Notice that the duality
test uses the flux \(F\), whose Euclidean \(p'\)-energy equals \(T_E\),
rather than the gradient \(\nabla w_E\) with an incorrect exponent. \(\square\)

\Needspace{6\baselineskip}\textbf{Corollary 4.2.} For \(q>0\) and \(\delta=\delta_p(H)\leq\eta\),

\[
\log R_q(H)\leq(-1+qC)\delta,
\qquad
\log R_q(H)\geq
\left(\frac q{p-1}-\frac{2pC}{n}\right)\delta.
\tag{4.13}
\]

\emph{Proof.} For the first estimate use \(\log(1-\delta)\leq-\delta\)
and \(\log(1+C\delta)\leq C\delta\). For the second use
\(\log(1-C\delta)\geq-2C\delta\), valid because \(C\delta\leq1/2\), and
\(-\log(1-\delta)\geq\delta\). Apply (4.6)-(4.7) to the logarithm of the
product. \(\square\)

\section{5. Persistence and quantitative
equivalence}\label{persistence-and-quantitative-equivalence}

With the constants from (2.11) and (4.5), choose

\[
q_-=
\min\left\{\frac{p-1}{2},\frac1{2C},
\frac{-\log(1-\eta)}{2\log M}\right\},
\tag{5.1}
\]

and

\[
q_+=(p-1)\max\left\{\frac{2pC}{n}+1,
\frac{\log(1/\ell)+1}{-\log(1-\eta)}\right\}.
\tag{5.2}
\]

These constants satisfy (1.4).

\emph{Proof of Theorem 1.1.} Suppose first that
\(0<\delta_p(H)\leq\eta\). Corollary 4.2 gives

\[
R_q(H)\leq e^{-\delta_p(H)/2}<1
\quad(0<q\leq q_-),
\tag{5.3}
\]

and

\[
R_q(H)\geq e^{\delta_p(H)}>1
\quad(q\geq q_+).
\tag{5.4}
\]

For \(\delta_p(H)\geq\eta\), Lemma 2.1 and (2.11) imply

\[
R_q(H)\leq(1-\eta)M^q
\leq\sqrt{1-\eta}<1
\quad(0<q\leq q_-),
\tag{5.5}
\]

whereas

\[
R_q(H)\geq\ell(1-\eta)^{-q/(p-1)}
\geq e>1
\quad(q\geq q_+).
\tag{5.6}
\]

These estimates include all degenerate seminorms. Finally,
\(\delta_p(H)=0\) forces \(H=1\) everywhere on \(S\) by continuity, so
\(H=E\). Equations (5.3)-(5.6) prove uniqueness and existence of the
stated optimizers directly.

For completeness, we establish both sides of (1.6). In the small-\(q\)
regime and for \(\delta=\delta_p(H)\leq\eta\), (5.3) gives

\[
D_q(H)\geq1-e^{-\delta/2}
\geq\tfrac12 e^{-1/2}\delta.
\tag{5.7}
\]

Since \(T_H\geq T_E\), the other side follows from (4.6):

\[
R_q(H)\geq(1-C\delta)^{p/n}
\geq1-\max\{1,p/n\}C\delta.
\tag{5.8}
\]

The last inequality follows from Bernoulli\textquotesingle s inequality
if \(p/n\geq1\) and from \((1-x)^{p/n}\geq1-x\) otherwise. In the
large-\(q\) regime, (5.4) implies \(D_q(H)\geq\delta\). For its upper
bound, \(\lambda_H\leq\lambda_E\) and (4.7) give

\[
R_q(H)\leq(1+C\delta)^q\leq1+C_q\delta,
\tag{5.9}
\]

where one may take \(C_q=qC\max\{1,(3/2)^{q-1}\}\).

It remains to treat \(\delta\geq\eta\). We always have \(\delta<1\). The
uniform gaps in (5.5)-(5.6) therefore give a positive constant times
\(\delta\) as a lower bound for \(D_q\). In the maximum regime
\(D_q\leq1\leq\delta/\eta\). In the minimum regime \(R_q\leq M^q\), so
\(D_q\leq M^q\delta/\eta\). Combining these constants with (5.7)-(5.9)
proves (1.6). \(\square\)

\section{6. Optimal supremum norm
stability}\label{optimal-supremum-norm-stability}

\Needspace{6\baselineskip}\textbf{Lemma 6.1 (angular concentration bound).} There is \(c_n>0\)
such that every \(H\in\mathcal H\) satisfies

\[
c_n\varepsilon(H)^{(n+1)/2}
\leq\delta_p(H)\leq p\varepsilon(H).
\tag{6.1}
\]

\emph{Proof.} If \(\varepsilon=0\), both assertions are immediate.
Otherwise choose \(\nu\in S\) with \(H(\nu)=1-\varepsilon\). The cap
calculation (3.2), which holds for every \(0<e\leq1\) without an inball
assumption, applies with \(e=\varepsilon\). On
\(C(\nu,\sqrt\varepsilon/8)\) we have \(1-H\geq\varepsilon/2\), and
hence \(1-H^p\geq\varepsilon/2\). Its measure is at least a dimensional
constant times \(\varepsilon^{(n-1)/2}\). This proves the lower bound.
The upper bound follows from \(1-t^p\leq p(1-t)\) for \(0\leq t\leq1\).
\(\square\)

\emph{Proof of Theorem 1.2.} The lower bound follows at once by
combining (1.6) and Lemma 6.1. For sharpness, take

\[
K_\varepsilon=
\overline B\cap\{x:|x\cdot e_1|\leq1-\varepsilon\},
\qquad H_\varepsilon=h_{K_\varepsilon},
\quad 0<\varepsilon<\tfrac12.
\tag{6.2}
\]

This centrally symmetric body contains \((1-\varepsilon)\overline B\)
and, since \(n\geq2\), touches \(S\) on the equator perpendicular to
\(e_1\). Consequently \(H_\varepsilon\) is a genuine normalized norm. At
\(e_1\) its support equals \(1-\varepsilon\), so
\(\varepsilon(H_\varepsilon)=\varepsilon\).

Whenever \(|\theta\cdot e_1|\leq1-\varepsilon\), the point \(\theta\)
itself belongs to \(K_\varepsilon\), and thus
\(H_\varepsilon(\theta)=1\). The angular deficit is supported on two
spherical caps of angular radius
\(\arccos(1-\varepsilon)\leq C\sqrt\varepsilon\). Their total measure is
at most \(C_n\varepsilon^{(n-1)/2}\). Also
\(1-H_\varepsilon^p\leq p\varepsilon\), whence

\[
\delta_p(H_\varepsilon)\leq C_{n,p}\varepsilon^{(n+1)/2}.
\tag{6.3}
\]

The upper bound in (1.6) proves (1.9). If a uniform estimate
\(D_q(H)\geq c\varepsilon(H)^\gamma\) held for some \(\gamma<(n+1)/2\),
evaluating it at \(H_\varepsilon\) and letting
\(\varepsilon\downarrow0\) would force a fixed positive \(c\) to be
zero. This proves optimality. \(\square\)

\section{7. Scope and further
questions}\label{scope-and-further-questions}

The argument identifies exact optimizers at finite exponents and gives a
common angular measure of their stability. The two regimes and the
nonlinear extension follow from the same comparison mechanism. None of
the constants in the construction is intended to be optimal.

The natural next question is to determine the full interval of \(q\) for
which \(E\) is extremal. For \(p=2\), the quadratic-seminorm conclusion
of {[}BFH, Theorem 4.10{]} reaches \(q=1\) in the maximizing problem;
extending that endpoint to arbitrary norms requires information not
supplied by (4.13). Another question is whether analogous persistence
holds on fixed domains other than balls. Our radial trial identities and
dual flux calculation depend on the ball geometry, and the present
theorem makes no assertion for general domains. Finally, constants
uniform as \(p\downarrow1\) or \(p\uparrow\infty\) are not obtained
here.

\sloppy\section{References}\label{references}

{[}AFTL{]} A. Alvino, V. Ferone, G. Trombetti and P.-L. Lions,
\emph{Convex symmetrization and applications}, Annales de
l\textquotesingle Institut Henri Poincaré C, Analyse non linéaire
\textbf{14} (1997), 275-293.
\href{https://www.numdam.org/item/AIHPC_1997__14_2_275_0/}{Primary
article and full text}.

{[}BFH{]} G. Buttazzo and R. Fernandes Horta, \emph{On the relations
between fundamental frequency and torsional rigidity in the case of
anisotropic energies}, arXiv:2603.09851v1 (2026).
\href{https://arxiv.org/html/2603.09851v1}{Versioned full text}.

{[}BHK{]} F. Besau, S. Hoehner and G. Kur, \emph{Intrinsic and dual
volume deviations of convex bodies and polytopes}, International
Mathematics Research Notices \textbf{2021} (22), 17456-17513.
\href{https://doi.org/10.1093/imrn/rnz277}{DOI};
\href{https://arxiv.org/abs/1905.08862}{author preprint}.

{[}HM{]} R. Fernandes Horta and M. Montenegro, \emph{Sharp
isoanisotropic estimates for fundamental frequencies of membranes and
connections with shapes}, Calculus of Variations and Partial
Differential Equations \textbf{64} (2025), article 273.
\href{https://doi.org/10.1007/s00526-025-03114-2}{Published article};
\href{https://arxiv.org/html/2406.18683v4}{arXiv:2406.18683v4}.

\end{document}
